\documentclass[11pt,letterpaper]{article}
\usepackage[T1]{fontenc}
\usepackage[utf8]{inputenc}
\usepackage{lmodern}
\usepackage[margin=1in,headheight=14pt]{geometry}
\usepackage{amsmath,amssymb,amsthm,mathtools,mathrsfs}
\usepackage{microtype,booktabs,array,enumitem}
\usepackage{xcolor}
\usepackage{fancyhdr}
\usepackage[unicode,hidelinks]{hyperref}
\usepackage[nameinlink,noabbrev]{cleveref}
\hypersetup{pdftitle={Exact Linearization Domains and Resonant Boundary Cycles in Surcomplex Dynamics},pdfauthor={Research manuscript prepared with ChatGPT},pdfsubject={Arbitrary-rank Hahn fields, polynomial dynamics, exact linearization, resonant cycles}}
\pagestyle{fancy}
\fancyhf{}
\fancyhead[L]{\small Surcomplex dynamics}
\fancyhead[R]{\small Exact domains and resonant cycles}
\fancyfoot[C]{\thepage}
\setlength{\parskip}{3pt}
\setlength{\emergencystretch}{2em}
\setlist{itemsep=3pt,topsep=5pt}
\numberwithin{equation}{section}
\theoremstyle{plain}
\newtheorem{theorem}{Theorem}[section]
\newtheorem{lemma}[theorem]{Lemma}
\newtheorem{proposition}[theorem]{Proposition}
\newtheorem{corollary}[theorem]{Corollary}
\theoremstyle{definition}
\newtheorem{definition}[theorem]{Definition}
\newtheorem{example}[theorem]{Example}
\theoremstyle{remark}
\newtheorem{remark}[theorem]{Remark}
\newcommand{\C}{\mathbb C}
\newcommand{\Q}{\mathbb Q}
\newcommand{\N}{\mathbb N}
\newcommand{\No}{\mathbf{No}}
\newcommand{\SC}{\mathbf{SC}}
\newcommand{\OO}{\mathcal O}
\newcommand{\mm}{\mathfrak m}
\newcommand{\supp}{\operatorname{supp}}
\newcommand{\res}{\operatorname{res}}
\newcommand{\id}{\operatorname{id}}
\newcommand{\ord}{\operatorname{ord}}
\newcommand{\lin}{\operatorname{lin}}
\newcommand{\sgn}{\operatorname{sgn}}
\newcommand{\coef}[2]{[#1^{#2}]}
\newcommand{\vv}{v}
\newcommand{\M}{\mathsf M}
\newcommand{\HS}{\mathscr B}
\newcommand{\smallo}{o_v}
\DeclareMathOperator{\Frac}{Frac}
\DeclareMathOperator{\Disc}{Disc}

\title{\textbf{Exact Linearization Domains and\newline Resonant Boundary Cycles\newline in Surcomplex Dynamics}\\[8pt]
\large Finite-return formulas, Hahn support certificates,\\
and a cancellation phase diagram}
\author{Research manuscript prepared with ChatGPT\\
\small For the Surreal project}
\date{21 September 2026}

\begin{document}
\maketitle
\thispagestyle{empty}
\begin{abstract}
Let $K=\C((t^\Gamma))$, where $\Gamma$ is any nonzero divisible ordered
abelian group, with no rank-one or discreteness assumption. For a polynomial
$f(z)=\lambda z+O(z^2)$ with $v(\lambda)=0$ and $\lambda$ not a root of unity,
we determine the exact maximal centered valuation ball on which the normalized
Schr\"oder conjugacy and its inverse are strongly Hahn-summable and isometric.
If the residue of $\lambda$ has order $q$, write
$f^{\circ q}(z)=\lambda^qz+\sum_{n\ge2}b_nz^n$ and
$\delta=v(\lambda^q-1)$. The ball is
\[
 \left\{z:v(z)>\max_{b_n\ne0}\frac{\delta-v(b_n)}{n-1}\right\}.
\]
Its bounding valuation sphere contains points of exact period $q$.
A finite ordinary polynomial determines their number with multiplicity,
their residue directions, and the first nonconstant term of their
multipliers. A separate support argument makes the conjugacy meaningful at
arbitrary valuation rank; a transfinite differential recursion further shows
that in the resonant-residue case all Hahn coefficient functions are
holomorphic on one common ordinary disk. This common-domain conclusion fails
in general for nonresonant residue, even though infinitesimal Hahn
linearization still exists. An explicit cubic family exhibits a transition
from one to two closest two-cycles and an exact improvement from a threshold
$\delta/2$ to $\delta/4$. All statements are stable under enlargement of the
Hahn workspace and therefore apply pointwise to $\No[i]$.

The exact resonant formula, its support-controlled and common-domain
realizations, and the resulting cancellation analysis are proposed as original
results. No named published conjecture is claimed solved, and priority is not
certified. The article supplies mathematical proofs and finite symbolic checks,
not a proof-assistant verification.
\end{abstract}

\newpage
{\small\setlength{\parskip}{0pt}\tableofcontents}
\newpage

\section{The question and the contribution}\label{sec:intro}

Local iteration asks a question not answered by locating the zeros of one
function: can a nonlinear map near a fixed point be changed into a linear map
by a single invertible analytic coordinate? The relevant equation is
\begin{equation}\label{eq:schroeder-intro}
 H(f(z))=\lambda H(z),\qquad H(0)=0,\quad H'(0)=1.
\end{equation}
Even when its formal solution is unique, one must still determine where the
series can be evaluated, where it is invertible, and what prevents extension
of the conjugacy.

Over surcomplex numbers, these are not merely ordinary convergence questions
with larger constants. At higher valuation rank, coefficient valuations alone
need not establish Hahn summability. Moreover, being Hahn-evaluable at every
infinitesimal argument is weaker than having ordinary holomorphic coefficient
functions on a common neighborhood. This article keeps those distinctions
throughout.

\subsection{Relation to the supplied repository}

The inspected catalogues of the \texttt{Surreal} repository describe extensive
work on Hahn-supported analysis, local analytic geometry, finite deformations,
residues, contours, divisors, polynomial algebra, and trigonometry
\cite{repo}. The present article develops polynomial iteration and normalized
linearization rather than restating those packages. The comparison was made
against the catalogues and a targeted repository search, not against every
proof in every source manuscript. In particular, the arguments below do not
assume the repository's more substantial unrefereed ring-theoretic theorems.
The prerequisites are the classical Hahn support lemma and algebraic closure
of the chosen Hahn field.

\subsection{What is proved}

The main results form a sequence rather than a collection of unrelated
observations.
\begin{enumerate}[label=\arabic*.]
\item \textbf{Exact linearization.} A finite maximum involving one return
iterate determines the maximal centered ball; the proof provides a uniform
support certificate for the conjugacy and its inverse
(\cref{thm:exact}).
\item \textbf{An algebraic obstruction to extension.} A residue polynomial
counts the closest resonant cycles and determines their leading multipliers
(\cref{thm:shell}).
\item \textbf{A stronger analytic realization.} In the resonant-residue case,
all Hahn coefficient functions of the scaled conjugacy are holomorphic on a
single ordinary disk (\cref{thm:coherent}). The proof uses a triangular family
of ordinary differential equations, not a convergence assertion about
successive surreal approximations.
\item \textbf{Cancellation and stability.} A two-parameter cubic family has
an exact, piecewise-linear threshold and an explicit change in the closest
cycle pattern. A finite-return perturbation certificate preserves the
threshold and matches simple cycles with a valuation bound
(\cref{thm:stability,thm:phase}).
\end{enumerate}

\subsection{Published antecedents and the novelty boundary}

Lindahl's equal-characteristic-zero theorem gives the exact polynomial
linearization disk when the residue multiplier is not a root of unity.
For residue order $q$, it gives a general lower radius
$|1-\lambda^q|^{1/q}/a$, where
$a=\max_{n\ge2}|a_n|^{1/(n-1)}$, together with an upper radius
\cite[Theorem~1.1]{lindahl-equal}. We do not claim that baseline
nonresonant result, or formal Schr\"oder linearization, as new. The resonant
formula here uses the actual coefficients of $f^{\circ q}$ and detects
cancellations discarded by a bound depending only on $a$, $q$, and
$|1-\lambda^q|$.

The classical support lemma \cite{neumann}, Hahn-field algebraic closure
\cite{poonen}, and ordinary one-variable differential equations are established
tools. What is offered for scrutiny is their use to obtain the exact
finite-return statement at arbitrary rank, the shell-cycle package, and the
common-domain lifting theorem. A targeted literature search did not locate
these statements in the form proved here. This is a proposed-originality
assessment, not an exhaustive priority claim. The results are not presented
as resolutions of an identified published open problem.

\section{Hahn workspaces and the analytic categories}\label{sec:setup}

\subsection{The field, its valuation, and balls}

Fix a nonzero divisible ordered abelian group $\Gamma$ and put
\[
 K=\C((t^\Gamma)).
\]
An element is a formal sum $a=\sum_{\gamma\in S}a_\gamma t^\gamma$ with
$S\subseteq\Gamma$ well ordered in the increasing order. Set
\[
 v(a)=\min\supp(a)\quad(a\ne0),\qquad v(0)=+\infty.
\]
Its valuation ring and maximal ideal are
\[
 \OO=\{a:v(a)\ge0\},\qquad \mm=\{a:v(a)>0\}.
\]
The residue field is canonically $\C$, and a bar denotes reduction modulo
$\mm$. In particular, every nonzero ordinary complex number, including every
nonzero integer, has valuation zero. The field $K$ is algebraically closed:
the relevant Hahn-field consequence of maximal completeness is recorded in
Poonen's Corollary~4 \cite{poonen}.

For $r\in\Gamma$ define
\begin{equation}\label{eq:balls}
 B_r=\{z\in K:v(z)>r\},\qquad S_r=\{z\in K:v(z)=r\}.
\end{equation}
Then $B_r=t^r\mm$. Increasing $r$ makes the ball \emph{smaller}.
We call $S_r$ the \emph{bounding shell} of $B_r$. It is not the topological
boundary: valuation balls are clopen. Nor is it a circle for the
ordered-surreal modulus.

A map $H:B_r\to B_r$ is an isometry when
$v(H(z)-H(w))=v(z-w)$. A linearization on a ball means an injective,
zero-fixing conjugacy with multiplication by $\lambda$ on a forward-invariant
centered valuation ball. The conjugacy constructed below is in fact a
bijection of that ball onto itself.

\subsection{Strong summability}

\begin{definition}\label{def:summable}
A set-indexed family $(u_i)_{i\in I}$ in $K$ is \emph{strongly Hahn-summable}
when $\bigcup_i\supp(u_i)$ is well ordered and each exponent belongs to the
support of only finitely many $u_i$. Its sum is formed coefficientwise.
\end{definition}

There is no limiting operation in this definition. All rearrangements below
are justified by this support condition. In particular, a statement such as
$\sum_{j\ge0}c^j=(1-c)^{-1}$, for $c\in\mm$, is a statement about a
summable family and multiplication, not about a convergent sequence of partial
sums in the full surreal fine topology.

\begin{lemma}[Positive-support lemma]\label{lem:neumann}
Let $S\subseteq\Gamma_{>0}$ be well ordered. Its additive monoid
\[
 \M=\langle S\rangle=\{s_1+\cdots+s_k:k\ge0,\ s_i\in S\}
\]
is well ordered. Each $\gamma\in\Gamma$ has only finitely many
representations by finite nondecreasing lists of elements of $S$, and hence
only finitely many representations by finite ordered lists.
\end{lemma}

This is the standard Neumann support lemma \cite{neumann}. A finite union of
well-ordered positive subsets may be used as $S$. The zero exponent belongs
to $\M$ through the empty list; zero is never a generator. The finite-list
conclusion, rather than an Archimedean estimate for list length, is what we
need at higher rank.

\subsection{A support algebra}

\begin{definition}\label{def:BM}
For $\M=\langle S\rangle$ as above, let
\[
 \HS_{\M}=
 \left\{\sum_{n\ge0}h_nX^n\in K[[X]]:
           \supp(h_n)\subseteq\M\text{ for every }n\right\}.
\]
Regrouping exponents identifies such a series with
$\sum_{\gamma\in\M}H_\gamma(X)t^\gamma$, where
$H_\gamma\in\C[[X]]$. This is a subring of
$\C[[X]]((t^\Gamma))$. Coefficients need not converge as ordinary complex
power series.
\end{definition}

The same notation is used for the evident nonnegative-support subring when
$S$ is empty, so $\M=\{0\}$. The uniform monoid is essential: an arbitrary
member of $\OO[[X]]$ need not belong to any such algebra.

\begin{lemma}[Evaluation, inversion, and isometry]\label{lem:evaluation}
Every $H\in\HS_\M$ is strongly Hahn-evaluable at every $u\in\mm$.
If $H=X+O(X^2)$, its formal inverse belongs to $\HS_\M$, and the evaluated
maps are mutually inverse isometric bijections of $\mm$. Formal composition
identities between zero-fixing series with integral unit linear coefficients
and a common positive-support certificate remain valid after evaluation.
\end{lemma}

\begin{proof}
For $u\ne0$, let $U=\supp(u)\subseteq\Gamma_{>0}$. The union of supports of
$h_nu^n$ lies in $\langle S\cup U\rangle$. Every occurrence of a monomial in
$u^n$ uses at least $n$ generators from $U$. By \cref{lem:neumann}, at a
fixed exponent only finitely many such lists, hence finitely many indices
$n$, occur. This proves both requirements in \cref{def:summable}.
The case $u=0$ is immediate.

The coefficients of the formal inverse of $X+\sum_{n\ge2}h_nX^n$ are
polynomials with integer coefficients in finitely many $h_n$. Thus their
supports also lie in $\M$. The substitutions in the inverse identities are
summable by the same list argument, so the evaluated functions are inverse.
For distinct $u,w\in\mm$, factor the divided difference:
\[
 \frac{H(u)-H(w)}{u-w}
 =1+\sum_{n\ge2}h_n\sum_{j=0}^{n-1}u^jw^{n-1-j}.
\]
Every nonconstant summand has positive valuation; the whole family is
summable by the same proof. The right side has residue one and valuation
zero. This proves the isometry assertion.

For compositions, expand both levels into products. At a fixed exponent the
number of positive generators is finite; at a fixed formal degree the number
of zero-exponent contributions is also finite, because the substituted
series has zero constant term. These two facts justify precisely the
rearrangements needed. A unit linear coefficient and its inverse may be
adjoined to the support certificate by separating their nonzero complex
constant terms and expanding the inverse in its positive part.
\end{proof}

\begin{example}[Why coefficient bounds alone are insufficient]\label{ex:bad-support}
Take $\Gamma=\Q+\Q\varepsilon$ with $0<\varepsilon<1/n$ for every positive
integer $n$. All coefficients of
\[
 A(X)=\sum_{n\ge2}t^{1/n}X^n
\]
have positive valuation. Nevertheless, at $u=t^\varepsilon$ the proposed
Taylor sum has exponents
\[
  \frac1n+n\varepsilon,
\]
which are strictly decreasing as $n$ increases. It is not a Hahn series.
Thus an assertion that ``all normalized coefficients are integral'' would
leave a real gap in an arbitrary-rank linearization proof. The union
$\{1/n:n\ge2\}$ is not well ordered, so this example does not satisfy the
uniform-support hypothesis of \cref{lem:evaluation}.
\end{example}

\subsection{Three function classes that must not be conflated}

The formal-coefficient ring $\C[[X]]((t^\Gamma))$ is distinct from the
radius-free convergent-germ ring $\C\{X\}((t^\Gamma))$. Both are distinct
from a common-domain ring
\[
 \OO(D)((t^\Gamma)),
\]
where every coefficient function is holomorphic on one fixed ordinary
complex disk $D$. A common-domain germ is obtained by allowing one subsequent
ordinary shrinking of $D$. These are the same distinctions emphasized in the
repository's analytic-geometry catalogue \cite{repo}.

The exact linearization theorem initially constructs a formal-coefficient
Hahn family that can be evaluated on the entire monad. Only
\cref{sec:coherent} establishes a common-domain representative, and it does
so only under the resonant-residue hypothesis.

\section{Two support-controlled linearization lemmas}\label{sec:linear-lemmas}

\subsection{The triangular recurrence}

Let
\[
 T(X)=\nu X+\sum_{j=2}^{D}d_jX^j,\qquad
 H(X)=X+\sum_{n\ge2}h_nX^n.
\]
When $\nu$ is nonzero and not a root of unity, comparison of degree $n$ in
$H(T(X))=\nu H(X)$ gives
\begin{equation}\label{eq:recurrence}
 (\nu-\nu^n)h_n
   =\sum_{j=1}^{n-1}h_j[X^n]T(X)^j,\qquad h_1=1.
\end{equation}
The denominators are nonzero, so this defines a unique normalized formal
conjugacy over any characteristic-zero field. It remains to control all of
its coefficients simultaneously.

\begin{lemma}[Nonresonant residue]\label{lem:nonresonant}
Suppose $T\in\OO[X]$, $v(\nu)=0$, and $\alpha=\bar\nu$ is not a root of
unity in $\C$. The normalized formal conjugacy and its inverse belong to
$\HS_\M$ for a positive-support monoid $\M$ depending only on the finitely
many coefficients of $T$.
\end{lemma}

\begin{proof}
Write $\nu=\alpha+e$, where $e\in\mm$. Let $S$ be the finite union of the
positive supports of $e$ and of the nonlinear coefficients $d_j$.
For every $n\ge2$, the expression $\nu-\nu^n$ has nonzero constant term
$\alpha-\alpha^n$. Its inverse is a formal power series in $e$ with ordinary
complex coefficients, so its support lies in $\langle\supp(e)\rangle$.
The recurrence \eqref{eq:recurrence} now proves inductively that every
$h_n$ is integral and supported in $\M=\langle S\rangle$.
The formal inverse has the same property by \cref{lem:evaluation}.
\end{proof}

There is no Diophantine lower bound on the ordinary complex numbers
$\alpha-\alpha^n$ in this proof. Those constants may be arbitrarily small
in the ordinary modulus and still all have Hahn valuation zero.

\subsection{Cancellation of a common infinitesimal divisor}

\begin{lemma}[Divisible displacement]\label{lem:displacement}
Let $0\ne c\in\mm$, $A\in X^2\OO[X]$, and
\begin{equation}\label{eq:T-displacement}
 T(X)=(1+c)X+cA(X)=X+cV(X),\qquad V=X+A.
\end{equation}
Then $1+c$ is not a root of unity. The normalized formal conjugacy of $T$
to $(1+c)X$, and its inverse, have all coefficients supported in one
positive-support monoid generated by the supports of $c$ and the positive
parts of the coefficients of $A$.
\end{lemma}

\begin{proof}
For a positive integer $n$, the first nonzero term of $(1+c)^n-1$ is $nc$;
all remaining terms have strictly larger valuation. Hence the multiplier is
not a root of unity.

For $n\ge2$ write
\begin{equation}\label{eq:denomcancel}
 (1+c)-(1+c)^n=-c(n-1)u_n(c),\qquad u_n(0)=1,
\end{equation}
where $u_n\in\Q[c]$. Whenever $j<n$, the coefficient
$[X^n]T(X)^j$ contains at least one nonlinear term, hence is divisible by
$c$ as a polynomial in $c$ and the coefficients of $A$. Substituting this
observation and \eqref{eq:denomcancel} into \eqref{eq:recurrence} cancels
one factor $c$ \emph{before} any valuation estimate is made.

The remaining division is by a nonzero ordinary integer and by $u_n(c)$.
The latter inverse is supported in the monoid generated by $\supp(c)$.
An induction on $n$ proves that every normalized coefficient has support in
one and the same monoid. The generating set is a finite union of well-ordered
positive supports and is therefore well ordered. Inversion of $H$ again uses
only polynomials in its coefficients.
\end{proof}

This cancellation is the basic reason the exact threshold is read from a
return map. Multiplying all small divisors in the recurrence would lose the
common factor present in every nonlinear numerator.

\subsection{An elementary root-counting lemma}

\begin{lemma}[Roots on the unit shell]\label{lem:root-count}
Let $Q\in\OO[X]$ with $Q(0)=1$, and suppose $P=\bar Q$ has degree
$M>0$. Then $Q$ has no root in $\mm$, and precisely $M$ roots of valuation
zero, counted with algebraic multiplicity. More precisely, for each nonzero
root $a\in\C$ of $P$, the roots of $Q$ with residue $a$ have total
multiplicity equal to the multiplicity of $a$ in $P$.
\end{lemma}

\begin{proof}
A root of positive valuation is impossible because the constant term would
strictly dominate all other terms. Factor in the algebraically closed field:
\[
 Q(X)=\prod_j(1-X/\beta_j),
\]
with roots repeated by multiplicity. Every $v(\beta_j)\le0$, so all factors
are integral. A factor with $v(\beta_j)<0$ reduces to $1$; a factor with
$v(\beta_j)=0$ reduces to $1-X/\bar\beta_j$. Reduction of this finite
factorization proves all assertions. In particular, a simple root of $P$
corresponds to exactly one simple root of $Q$ in that residue class.
\end{proof}

\section{The exact maximal ball}\label{sec:exact}

Let
\begin{equation}\label{eq:f-main}
 f(z)=\lambda z+\sum_{n=2}^{d}a_nz^n\in K[z],\qquad d\ge2,
 \quad a_d\ne0,\quad v(\lambda)=0,
\end{equation}
and assume throughout that $\lambda$ is not a root of unity. Put
\begin{equation}\label{eq:eta}
 \eta=\eta(f)=\max_{a_n\ne0}\frac{-v(a_n)}{n-1},
 \qquad \alpha=\bar\lambda\in\C^\times.
\end{equation}
All maxima in this article are finite maxima over nonzero polynomial
coefficients. Rational division makes sense because $\Gamma$ is divisible.

\begin{theorem}[Exact finite-return linearization]\label{thm:exact}
Define $r=r(f)$ as follows.

\emph{Nonresonant residue.} If $\alpha$ is not a root of unity, set
\begin{equation}\label{eq:nonres-r}
 r=\eta(f).
\end{equation}

\emph{Resonant residue.} If $\alpha$ has exact order $q\ge1$, write
\begin{equation}\label{eq:g-main}
 g(z)=f^{\circ q}(z)=\lambda^qz+\sum_{n=2}^{d^q}b_nz^n,
 \qquad c=\lambda^q-1,\qquad \delta=v(c)>0,
\end{equation}
and set
\begin{equation}\label{eq:exact-r}
 \boxed{\displaystyle
 r=\max_{b_n\ne0}\frac{\delta-v(b_n)}{n-1}.}
\end{equation}
In this case $r>\eta(f)$.

Let $H(z)=z+\sum_{n\ge2}h_nz^n$ be the unique normalized formal solution
of $H\circ f=\lambda H$. Then $H$ and its formal inverse are strongly
Hahn-evaluable on $B_r$, and their evaluated maps are mutually inverse
isometries of $B_r$. In particular, $f:B_r\to B_r$ is conjugate to
multiplication by $\lambda$.

More strongly, if $\rho=t^r$, then
\begin{equation}\label{eq:H-scaled}
 H_r(X)=\rho^{-1}H(\rho X)
       =X+\sum_{n\ge2}h_n\rho^{n-1}X^n
\end{equation}
and its inverse have a uniform positive-support certificate. Consequently
\begin{equation}\label{eq:hn-bound}
 v(h_n)+(n-1)r\ge0\qquad(n\ge2),
\end{equation}
but the certificate is stronger than this inequality.

No strictly larger centered valuation ball admits an injective conjugacy of
$f$ to multiplication by $\lambda$. In the nonresonant case the obstruction
is a nonzero $p\in S_r$ with $f(p)=0$. In the resonant case it is a nonzero
$p\in S_r$ of exact period $q$.
\end{theorem}

\subsection{Proof in the nonresonant case}

Take $\rho=t^\eta$ and define $F(X)=\rho^{-1}f(\rho X)$. Its nonlinear
coefficients have valuations
\[
 v(a_n)+(n-1)\eta\ge0.
\]
At least one has valuation zero. Hence $F\in\OO[X]$ and its reduction is
nonlinear. By \cref{lem:nonresonant,lem:evaluation}, its normalized
conjugacy and inverse are strongly summable, mutually inverse isometries on
$\mm$. Scaling back proves all existence and support assertions.

To see maximality, apply \cref{lem:root-count} to
\[
 Q(X)=\frac{F(X)}{\lambda X}
     =1+\sum_{n\ge2}\frac{a_n\rho^{n-1}}{\lambda}X^{n-1}.
\]
Its reduction is nonconstant. Thus $Q$ has a root $u$ of valuation zero,
and $p=\rho u\in S_\eta$ satisfies $f(p)=0=f(0)$. Any centered valuation
ball strictly containing $B_\eta$ contains this entire shell. An injective
conjugacy on that ball would imply
$\lambda H(p)=H(f(p))=H(f(0))=\lambda H(0)$, a contradiction.

\subsection{Why the resonant threshold lies inside the injectivity scale}

Suppose $\alpha$ has order $q$. The multiplier $\lambda^q$ has residue one,
and it is not one because $\lambda$ is not a root of unity. Thus $c\ne0$
and $\delta>0$.

Scale initially by $\sigma=t^\eta$. The polynomial
$F_\eta(X)=\sigma^{-1}f(\sigma X)$ is integral and its reduction has degree
at least two. Reduction commutes with polynomial iteration, so
\[
 \overline{\sigma^{-1}g(\sigma X)}
       =\bar F_\eta^{\circ q}(X)
\]
also has degree at least two. Therefore some $n\ge2$ satisfies
\[
 v(b_n)+(n-1)\eta=0.
\]
For this index,
\[
 \frac{\delta-v(b_n)}{n-1}
    =\eta+\frac{\delta}{n-1}>\eta.
\]
Taking the maximum proves $r>\eta$. In particular, at the final scale
$\rho=t^r$, the polynomial
\begin{equation}\label{eq:F-scaled}
 F(X)=\rho^{-1}f(\rho X)
\end{equation}
has reduction exactly $\alpha X$: every nonlinear coefficient has positive
valuation. It maps $\mm$ to itself and preserves the residue direction of a
unit up to multiplication by $\alpha$.

\subsection{Linearizing the return and recovering the original map}

At the scale $\rho=t^r$, write
\begin{align}
 G(X)&=\rho^{-1}g(\rho X)=(1+c)X+cA(X),\label{eq:G}\\
 A(X)&=\sum_{n\ge2}\frac{b_n\rho^{n-1}}{c}X^n\in X^2\OO[X].\label{eq:A}
\end{align}
Integrality follows exactly from \eqref{eq:exact-r}. The divisible-displacement
lemma constructs the normalized conjugacy $J$ of $G$ to $(1+c)X$, with a
uniform support certificate.

It remains to verify that $J$ linearizes $F$, not merely its $q$th iterate.
In the formal power-series group, the map
\[
 T=J\circ F\circ J^{-1}
\]
commutes with multiplication by $\mu=1+c$, because $F$ commutes with
$G=F^{\circ q}$. If $T(X)=\sum_{n\ge1}\tau_nX^n$, this commutation says
\[
 \tau_n(\mu^n-\mu)=0.
\]
The multiplier $\mu$ is not a root of unity, so all $\tau_n$ with $n\ge2$
vanish. Its linear coefficient is $\lambda$, hence $T(X)=\lambda X$.
Thus $J$ is the unique normalized conjugacy of $F$.

Adjoin the positive supports of the finitely many coefficients of $F$ to the
support certificate when evaluating this identity. All relevant compositions
are legitimate by \cref{lem:evaluation}. Scaling by $\rho$ proves the
existence, isometry, and coefficient assertions of \cref{thm:exact}.

\subsection{Exact-period obstruction}

Define the integral polynomial
\begin{equation}\label{eq:Q-shell}
 Q(X)=\frac{G(X)-X}{cX}
     =1+\sum_{n\ge2}\frac{b_n\rho^{n-1}}{c}X^{n-1}.
\end{equation}
At least one nonlinear coefficient of $G$ attains equality in the maximum,
so $P=\bar Q$ is nonconstant. By \cref{lem:root-count}, $Q$ has a unit
root $u$. Then $p=\rho u$ lies on $S_r$ and $g(p)=p$.

The least period $k$ of $p$ divides $q$. Every point in this finite orbit
lies on $S_r$, and reduction of the normalized orbit multiplies by
$\alpha$ at each step. If $a=\bar u\ne0$, then
$\alpha^ka=a$, so $q$ divides $k$. Thus $k=q$.

On any larger centered valuation ball, an injective conjugacy would give
\[
 H(p)=H(f^{\circ q}(p))=\lambda^qH(p).
\]
Since $\lambda^q\ne1$, this forces $H(p)=0=H(0)$, a contradiction. This
completes the proof of \cref{thm:exact}.\hfill$\square$

\begin{remark}[What is exact]\label{rem:exact-scope}
The conclusion determines the maximal \emph{centered valuation ball}, not the
full pointwise summability locus of the formal series, and not a maximal
non-ball invariant domain. The series can have exceptional evaluations on
its bounding shell. None can provide an injective conjugacy on a ball that
includes a forbidden cycle or collision.
\end{remark}

\section{The shell polynomial and the closest cycles}\label{sec:shell}

Continue in the resonant-residue case and retain $r$, $\rho$, $c$, $Q$ from
\cref{sec:exact}. Write
\begin{equation}\label{eq:P}
 P(X)=\bar Q(X),\qquad
 M=\deg P
   =\max\left\{n-1:v(b_n)+(n-1)r=\delta\right\}.
\end{equation}
This is a finite ordinary complex polynomial with $P(0)=1$. It is the
initial polynomial on the first face of the return equation seen from the
origin.

\begin{theorem}[Resonant shell cycles]\label{thm:shell}
The following assertions hold.
\begin{enumerate}[label=(\roman*)]
\item The equation $f^{\circ q}(z)=z$ has exactly $M$ nonzero solutions on
$S_r$, counted with algebraic multiplicity. Every such solution has exact
period $q$, and there are no nonzero periodic points in $B_r$.
\item $P(\alpha X)=P(X)$; equivalently, $P\in\C[X^q]$. In particular,
$q$ divides $M$.
\item If $P$ is squarefree, there is one simple point
$p_a=\rho(a+u_a)$, $u_a\in\mm$, for every root $a$ of $P$. These points
form exactly $M/q$ distinct cycles, and $f(p_a)=p_{\alpha a}$.
\item For each such simple point,
\begin{equation}\label{eq:multiplier}
 \res\left(\frac{(f^{\circ q})'(p_a)-1}{c}\right)=aP'(a),
 \qquad v\bigl((f^{\circ q})'(p_a)-1\bigr)=\delta.
\end{equation}
\end{enumerate}
\end{theorem}

\begin{proof}
The counting and residue-class assertions follow from
\cref{lem:root-count} applied to $Q$. Exact period $q$ was proved in
\cref{sec:exact}. A nonzero periodic point in $B_r$ would be sent by the
injective conjugacy to a nonzero periodic point of multiplication by
$\lambda$, which has none.

For the symmetry, let $F=\rho^{-1}f(\rho X)$ and write
\[
 G=F^{\circ q}=X+cW(X),\qquad W=XQ.
\]
The identity $F\circ G=G\circ F$ gives
\[
 \frac{F(X+cW(X))-F(X)}c=W(F(X)).
\]
This is an identity of integral polynomials. Reduction leaves only the
first Taylor term on the left, since $c\in\mm$. As $\bar F=\alpha X$,
\[
 \alpha\bar W(X)=\bar W(\alpha X).
\]
Substitution of $\bar W=XP$ yields $P(\alpha X)=P(X)$. Since $\alpha$
has exact order $q$, only degrees divisible by $q$ occur in $P$.

When $P$ is squarefree, the uniqueness in each residue class from
\cref{lem:root-count} implies $f(p_a)=p_{\alpha a}$. The action of
$\alpha$ on nonzero complex numbers has orbits of length $q$, so the
number of cycles is $M/q$.

Finally, differentiation of $G=X+cW$ gives
$G'(X)=1+cW'(X)$. At $p_a=\rho u$, $\bar u=a$, the derivative of the
unscaled iterate equals $G'(u)$. Reduction gives
\[
 \overline{W'(u)}=P(a)+aP'(a)=aP'(a)\ne0.
\]
This proves both assertions in \eqref{eq:multiplier}.
\end{proof}

\begin{remark}[Multiplicity is not the number of distinct points]
If $P$ has a multiple root, the theorem still gives a multiplicity count.
It does not say that several actual periodic points have collided: higher
Hahn coefficients can split one repeated initial root into distinct points
with the same residue direction. In that situation the equality
$v((f^{\circ q})'(p)-1)=\delta$ need not hold, because the displayed residue
may vanish.
\end{remark}

\begin{corollary}[A finite obstruction certificate]\label{cor:finite-certificate}
In the resonant case the data $q$, $c$, $r$, and $P$ certify exact maximality
without calculating a single coefficient of the infinite conjugacy. At most
$d^q-1$ shell points occur with multiplicity, and their count is a positive
multiple of $q$.
\end{corollary}

\begin{proof}
The return polynomial has degree $d^q$, and $M=\deg P\le d^q-1$.
The exactness and positivity follow from \cref{thm:exact,thm:shell}.
\end{proof}

\section{A common-domain lifting theorem}\label{sec:coherent}

The preceding support arguments establish more than formal existence, but
less than ordinary analyticity of every coefficient function. In the resonant
case a further argument supplies that missing regularity.

\begin{theorem}[Coherent realization of resonant linearization]\label{thm:coherent}
Under the resonant-residue hypotheses of \cref{thm:exact}, the scaled
conjugacy $H_r$ has an expansion
\begin{equation}\label{eq:coherent-expansion}
 H_r(X)=\sum_{\gamma\in\M}H_\gamma(X)t^\gamma
       \quad\text{in }\OO(D)((t^\Gamma))
\end{equation}
for a single ordinary complex disk $D$ centered at zero. Here $\M$ is a
positive-support monoid, $H_0(0)=0$, $H_0'(0)=1$, and
$H_\gamma\in X^2\OO(D)$ for $\gamma>0$.

One may take any sufficiently small ordinary disk on which the shell
polynomial $P$ has no zero. Its leading coefficient is
\begin{equation}\label{eq:H0}
 \boxed{\displaystyle
 H_0(X)=X\exp\left(\int_0^X
       \left(\frac1{sP(s)}-\frac1s\right)\,ds\right).}
\end{equation}
The integrand is holomorphic at zero. After an ordinary shrinking, the formal
inverse of $H_r$ also has all Hahn coefficient functions holomorphic on one
common ordinary neighborhood.
\end{theorem}

\begin{proof}
Use $G=X+cV$ from \eqref{eq:G}--\eqref{eq:A}, so $V=X+A=XQ$ and
\begin{equation}\label{eq:V0}
 V_0:=\bar V=XP(X).
\end{equation}
All coefficients of the normalized conjugacy lie in a common monoid $\M$
by \cref{lem:displacement}. Taylor expansion of the formal identity
$H_r(X+cV)=(1+c)H_r(X)$, followed by cancellation of $c$, gives
\begin{equation}\label{eq:differential-schroeder}
 V H_r'+\sum_{k\ge2}\frac{c^{k-1}}{k!}V^kH_r^{(k)}=H_r.
\end{equation}
This identity is first interpreted coefficientwise in $X$. At each degree,
only finitely many terms contribute because $V$ has order one. It may also
be regrouped at a fixed Hahn exponent: each term with index $k$ uses at
least $k-1$ positive-support generators from $c$, so
\cref{lem:neumann} makes that regrouping finite.

At exponent zero, \eqref{eq:differential-schroeder} becomes
\begin{equation}\label{eq:ode0}
 V_0H_0'=H_0.
\end{equation}
Choose an ordinary disk $D$ centered at zero on which $P$ has no zero.
Because $P(0)=1$, the function $1/(XP(X))-1/X$ has a removable
singularity at zero. Formula \eqref{eq:H0} defines the unique solution of
\eqref{eq:ode0} with $H_0'(0)=1$. The quotient $H_0/X$ is nonvanishing
on $D$.

Now let $\gamma>0$ in the well-ordered monoid $\M$, and assume
$H_\beta$ has been constructed on this same disk for every $\beta<\gamma$.
At exponent $\gamma$, isolate the terms containing $H_\gamma$ in
\eqref{eq:differential-schroeder}. They are precisely
\begin{equation}\label{eq:ode-gamma}
 L H_\gamma:=V_0H_\gamma'-H_\gamma=R_\gamma,
\end{equation}
where
\begin{equation}\label{eq:R-gamma}
 R_\gamma=-[t^\gamma]\left(
 (V-V_0)H_r'+
 \sum_{k\ge2}\frac{c^{k-1}}{k!}V^kH_r^{(k)}\right).
\end{equation}
Every appearance of a coefficient of $H_r$ on the right is multiplied by a
positive exponent. Thus only $H_\beta$ with $\beta<\gamma$ occurs. The
support lemma makes this a \emph{finite} expression in earlier coefficient
functions and their derivatives. Consequently it is holomorphic on the
same $D$; no new shrinking is required.

Moreover, $V-V_0$ vanishes to order at least two, since the coefficient of
$X$ in $V$ is exactly one. The terms with $k\ge2$ also vanish to order at
least two. Hence $R_\gamma\in X^2\OO(D)$. Variation of the constant in
\eqref{eq:ode0} gives the unique normalized solution
\begin{equation}\label{eq:gamma-integral}
 H_\gamma(X)=H_0(X)\int_0^X
       \frac{R_\gamma(s)}{V_0(s)H_0(s)}\,ds.
\end{equation}
The apparent singularity at zero is removable: the numerator vanishes to
order at least two and the denominator has order exactly two. Elsewhere on
$D$ the denominator is nonzero. Thus $H_\gamma\in X^2\OO(D)$.

Transfinite induction on $\M$ constructs all coefficients on one fixed disk.
They have the same formal Taylor expansions as the unique solution of the
Schr\"oder recurrence. Therefore this analytic family is the regrouping of
$H_r$ already constructed, not a different conjugacy.

For the inverse, first shrink to ordinary neighborhoods on which $H_0$ is
biholomorphic. The zero-exponent coefficient of the inverse is $H_0^{-1}$.
At any positive exponent, the coefficient equation for composition is linear
in the unknown inverse coefficient, with coefficient
$H_0'(H_0^{-1}(Y))$, which never vanishes on a sufficiently small fixed
neighborhood. All other terms depend on earlier exponents through finitely
many ordinary derivatives and products. The same support induction produces
a common-domain inverse, again without repeated shrinking.
\end{proof}

\subsection{Why the differential recursion is genuinely triangular}

It would be incorrect to justify \eqref{eq:R-gamma} by saying that
$(k-1)v(c)\to+\infty$: this can fail to be cofinal in an arbitrary-rank
value group. The correct statement is that, at a \emph{specified} exponent
$\gamma$, infinitely many decompositions into generators from one
well-ordered positive support are impossible. This simultaneously controls
the Taylor index, the Hahn products, and the number of earlier coefficients
entering the right side. The construction is transfinite but each individual
coefficient equation has finite dependence.

The common-domain theorem asserts no uniform ordinary norm bound on the
family $(H_\gamma)$. Coefficient functions may become arbitrarily large or
have arbitrarily large derivatives. What remains uniform is their ordinary
domain of holomorphy and the Hahn support certificate.

\subsection{Two analytic structures, one finite polynomial}

The same shell polynomial has two different roles. Its nonzero roots give the
closest periodic directions in the valued field, while its absence of zeros
on a small ordinary disk permits the coefficientwise differential equation
to be solved there. Formula \eqref{eq:H0} provides the precise bridge.
This does not identify the surcomplex ball $B_r$ with the ordinary disk $D$:
all points of the normalized monad have ordinary residue zero, whereas $D$
also contains many nonzero ordinary points. The latter is a domain for the
coefficient functions, not a claim of a larger centered surcomplex
linearization ball.

\section{A finite-return perturbation certificate}\label{sec:stability}

An exact threshold is most useful when one can tell which changes preserve
it. The following certificate refers to the finite return coefficients, not
to the infinite linearizing series.

\begin{theorem}[Stability above the first return face]\label{thm:stability}
Let $f$ satisfy the resonant hypotheses of \cref{thm:exact}, with multiplier
$\lambda$, return coefficients $b_n$, and data $q,c,\delta,r,P$.
Let $\widetilde f$ be another nonlinear polynomial with the same multiplier
$\lambda$, and write
\[
 \widetilde f^{\circ q}(z)=\lambda^qz+\sum_{n\ge2}\widetilde b_nz^n.
\]
Pad either coefficient list with zeros as needed. Suppose
\begin{equation}\label{eq:stable-assumption}
 v(\widetilde b_n-b_n)+(n-1)r>\delta
 \quad\text{for every nonzero difference.}
\end{equation}
Then $r(\widetilde f)=r(f)$ and the shell polynomial of $\widetilde f$
is also $P$.

If $P$ is squarefree and at least one coefficient changes, put
\begin{equation}\label{eq:epsilon}
 \epsilon=\min_{\widetilde b_n\ne b_n}
     \bigl(v(\widetilde b_n-b_n)+(n-1)r-\delta\bigr)>0.
\end{equation}
The closest cycles have a canonical matching by residue direction. Matched
points satisfy
\begin{equation}\label{eq:stable-distance}
 v(\widetilde p_a-p_a)\ge r+\epsilon.
\end{equation}
The matching intertwines $f$ and $\widetilde f$ on these finite cycle sets,
and their leading multiplier ratios in \eqref{eq:multiplier} agree.
\end{theorem}

\begin{proof}
At $\rho=t^r$ the two normalized displacement polynomials are
\[
 Q=1+\sum_{n\ge2}\frac{b_n\rho^{n-1}}cX^{n-1},\qquad
 \widetilde Q=1+\sum_{n\ge2}
                \frac{\widetilde b_n\rho^{n-1}}cX^{n-1}.
\]
The hypothesis says that every coefficient of $\widetilde Q-Q$ has positive
valuation. Thus $\widetilde Q$ is integral with the same nonconstant
reduction $P$. The inequalities defining the maximum hold at $r$ for all
$\widetilde b_n$, and at least one remains an equality. The exact formula
therefore gives $r(\widetilde f)=r$.

For a simple root $a$ of $P$, let $u=p_a/\rho$ and
$\widetilde u=\widetilde p_a/\rho$. Both have residue $a$.
The divided difference
\[
 \frac{Q(u)-Q(\widetilde u)}{u-\widetilde u}
\]
when $u\ne\widetilde u$ is integral with nonzero residue $P'(a)$.
Moreover,
\[
 Q(u)=0,\qquad
 Q(\widetilde u)=(Q-\widetilde Q)(\widetilde u),
\]
and the latter has valuation at least $\epsilon$. Hence
$v(u-\widetilde u)\ge\epsilon$, which is \eqref{eq:stable-distance}.
If the two roots coincide the inequality is automatic. The direction rule
$a\mapsto\alpha a$ and the multiplier formula from \cref{thm:shell}
prove the remaining assertions.
\end{proof}

The certificate is sufficient, not necessary. A perturbation may change the
first return face and nevertheless preserve its maximal slope. Also,
\eqref{eq:stable-assumption} is a condition on the \emph{return} polynomial;
one should not silently replace it by the same inequalities on the original
coefficients of $f$.

\section{Sharp examples and a cancellation phase diagram}\label{sec:examples}

\subsection{A tangent binomial}

Let
\begin{equation}\label{eq:binomial}
 f(z)=(1+c)z+az^m,\qquad 0\ne c\in\mm,\quad a\ne0,\quad m\ge2.
\end{equation}
Here $q=1$ and
\[
 r=\frac{v(c)-v(a)}{m-1}.
\]
The nonzero fixed points are exactly
\[
 p^{m-1}=-c/a.
\]
All $m-1$ are simple and lie on $S_r$. At each of them the multiplier is
known exactly, not just to leading order:
\begin{equation}\label{eq:binomial-mult}
 f'(p)=1+c+ma p^{m-1}=1-(m-1)c.
\end{equation}
If $A=\res(a\rho^{m-1}/c)$, the shell polynomial is
$P(X)=1+AX^{m-1}$ and
\[
 H_0(X)=X(1+AX^{m-1})^{-1/(m-1)},
\]
using the ordinary branch normalized to one at zero. This already shows why
a degree-independent small-divisor loss can be wasteful.

\subsection{A cubic with two suppressed return terms}\label{subsec:tuned}

Let $\delta>0$, set $e=t^\delta$, and consider
\begin{equation}\label{eq:tuned-cubic}
 f_e(z)=(-1+e)z+z^2-z^3.
\end{equation}
Its residue multiplier has order two, $\eta(f_e)=0$, and
$c=\lambda^2-1=e(e-2)$ has valuation $\delta$. The nonlinearity of the
second iterate is not guessed from the original quadratic term. Its exact
coefficients are:
\begin{center}
\begin{tabular}{c l}
\toprule
Degree $n$ & Coefficient $b_n$ in $f_e^{\circ2}$\\
\midrule
2 & $e(e-1)$\\
3 & $-e(e-2)(e-1)$\\
4 & $-e(3e-4)$\\
5 & $3e^2-9e+4$\\
6 & $6(e-1)$\\
7 & $-3(e-2)$\\
8 & $-3$\\
9 & $1$\\
\bottomrule
\end{tabular}
\end{center}
At $e=0$ the return displacement factors as
\begin{equation}\label{eq:tuned-factor}
 f_0^{\circ2}(z)-z
   =z^5(z^2-2z+2)(z^2-z+2).
\end{equation}
The cubic and quartic terms have both disappeared. Their actual valuations,
not the separate sizes of the terms that cancelled, enter \eqref{eq:exact-r}.
The degree-five coefficient is a unit, and the exact maximum is
\begin{equation}\label{eq:tuned-r}
 \boxed{r(f_e)=\delta/4.}
\end{equation}
For $\rho=t^{\delta/4}$,
\[
 P(X)=1-2X^4.
\]
There are four simple closest period-two points, in two cycles. Their
residue directions satisfy $a^4=1/2$. Since $aP'(a)=-4$,
\begin{equation}\label{eq:tuned-mult}
 (f_e^{\circ2})'(p_a)=1-4c+\smallo(c)
                     =1+8e+\smallo(e),
\end{equation}
where $R=\smallo(c)$ means $v(R)>v(c)$.
The common-domain leading conjugacy is
\[
 H_0(X)=X(1-2X^4)^{-1/4}.
\]

For comparison, the untuned quadratic $(-1+e)z+z^2$ has threshold
$\delta/2$ and shell polynomial $1+X^2$. Thus the cubic in
\eqref{eq:tuned-cubic} admits an exactly larger linearization ball even
though it has the same multiplier and the same basic coefficient scale
$\eta=0$. The difference is produced by cancellation in a finite return.

\subsection{The two-parameter transition}

The cancellation can be broken on a second valuation scale. Let
\begin{equation}\label{eq:phase-family}
 f_{e,s}(z)=(-1+e)z+z^2+(-1+s)z^3,
 \qquad e=t^\delta,\quad s=t^\sigma,
 \quad\delta,\sigma>0.
\end{equation}
Both parameters are monomials with coefficient one. No Archimedean
relationship between their positive exponents is assumed.

\begin{theorem}[Cancellation phase diagram]\label{thm:phase}
For \eqref{eq:phase-family},
\begin{equation}\label{eq:phase-r}
 \boxed{\displaystyle
 r(f_{e,s})=\max\left\{\frac{\delta-\sigma}{2},
                         \frac\delta4\right\}.}
\end{equation}
The shell polynomial, cycle count, and leading return-multiplier ratios are
as follows.
\begin{center}
\small
\begin{tabular}{@{}llll@{}}
\toprule
Regime & $P(X)$ & Two-cycles & $aP'(a)$\\
\midrule
$\sigma<\delta/2$ & $1+X^2$ & one & $-2$\\
$\sigma=\delta/2$ & $1+X^2-2X^4$ & two & $-6$, $-3$\\
$\sigma>\delta/2$ & $1-2X^4$ & two & $-4$\\
\bottomrule
\end{tabular}
\end{center}
Every listed initial root is simple. In the middle row the ratio $-6$
occurs for $a^2=1$, and $-3$ for $a^2=-1/2$.
\end{theorem}

\begin{proof}
For the general cubic $f(z)=(-1+e)z+z^2+bz^3$, direct composition gives
\begin{align}
 b_2&=e(e-1),\nonumber\\
 b_3&=(e-1)(be^2-2be+2b+2),\nonumber\\
 b_4&=3be^2-4be+b+1,\nonumber\\
 b_5&=b(3be^2-6be+3b+3e-1),\label{eq:general-cubic}\\
 b_6&=b(6be-5b+1),\nonumber\\
 b_7&=3b^2(be-b+1),\nonumber\\
 b_8&=3b^3,\qquad b_9=b^4.\nonumber
\end{align}
Substitute $b=-1+s$. The degree-three coefficient has leading terms
$-2e-2s$; the degree-four coefficient has leading terms $4e+s$.
Because $e$ and $s$ have leading coefficient one, no cancellation occurs
between those pairs even when $\delta=\sigma$. Thus
\[
 v(b_2)=\delta,\qquad
 v(b_3)=v(b_4)=\min\{\delta,\sigma\},\qquad
 v(b_n)=0\quad(5\le n\le9).
\]
The maximum in \eqref{eq:exact-r} is therefore the larger of
$(\delta-\min\{\delta,\sigma\})/2$ and $\delta/4$, which is
\eqref{eq:phase-r}.

If $\sigma<\delta/2$, take $r=(\delta-\sigma)/2$. Only the
normalized degree-three return coefficient has nonzero residue in $Q-1$,
and that residue is one. Thus $P=1+X^2$.
If $\sigma>\delta/2$, take $r=\delta/4$. Only degree five survives,
with residue $4/(-2)=-2$, giving $P=1-2X^4$.
At equality both terms survive, giving
\[
 P=1+X^2-2X^4=(1-X^2)(1+2X^2).
\]
All three polynomials are squarefree. The cycle counts follow from
\cref{thm:shell}. Finally, reduction of $XP'(X)$ at their roots gives
respectively $-2$, $-4$, and, in the middle case,
\[
 2a^2-8a^4=
 \begin{cases}-6,&a^2=1,\\-3,&a^2=-1/2.\end{cases}
\]
\end{proof}

The leading coefficient of the coherent conjugacy can also be read off in
all three regimes:
\begin{equation}\label{eq:phase-H0}
 H_0(X)=
 \begin{cases}
 X(1+X^2)^{-1/2},&\sigma<\delta/2,\\[2pt]
 X(1-X^2)^{-1/6}(1+2X^2)^{-1/3},&\sigma=\delta/2,\\[2pt]
 X(1-2X^4)^{-1/4},&\sigma>\delta/2.
 \end{cases}
\end{equation}
These formulas use branches equal to one at zero. Differentiating their
logarithms verifies $XP(X)H_0'(X)=H_0(X)$ directly.

\subsection{A repeated initial root at the transition}

Replace $s$ in the equality regime by $\kappa t^{\delta/2}$, where
$\kappa\in\C^\times$. The same calculation gives
\[
 P_\kappa(X)=1+\kappa X^2-2X^4.
\]
As a polynomial in $Y=X^2$, its discriminant is
\[
 \Disc_Y(1+\kappa Y-2Y^2)=\kappa^2+8.
\]
For $\kappa^2+8\ne0$, the four shell points are simple in distinct residue
directions. For $\kappa^2+8=0$, the initial polynomial has repeated roots.
This is an initial-direction degeneracy, not by itself an assertion of an
actual multiple periodic point. Determining the finer splitting there
requires higher Hahn terms of the return equation. The exact threshold
remains $\delta/4$, since its degree-four term in $P_\kappa$ cannot vanish.

\subsection{A genuinely higher-rank example}\label{subsec:rank2}

Let $\Gamma=\Q+\Q\Omega$ with $\Omega>n$ for every ordinary positive
integer $n$. Consider
\begin{equation}\label{eq:rank2}
 f(z)=(1+t^\Omega)z+t^3z^2+z^4.
\end{equation}
Here $q=1$, $\eta=0$, and
\[
 r=\max\{\Omega-3,\Omega/3\}=\Omega-3.
\]
At $\rho=t^{\Omega-3}$, the shell polynomial is $P(X)=1+X$. There is
one closest nonzero fixed point
\[
 p=\rho(-1+u),\qquad u\in\mm,
\]
and its multiplier satisfies $f'(p)=1-t^\Omega+\smallo(t^\Omega)$.

The other two nonzero fixed points are much farther from zero. Indeed,
all three solve $z^3+t^3z+t^\Omega=0$. Factoring off $p$ gives
\[
 z^3+t^3z+t^\Omega
  =(z-p)\bigl(z^2+pz+p^2+t^3\bigr).
\]
Scale the quadratic factor by $z=t^{3/2}Y$ and divide by $t^3$.
Its reduction is $Y^2+1$, because $v(p)=\Omega-3$ dominates every
ordinary rational. The two other roots therefore have valuation $3/2$ and
normalized residue directions $\pm i$.

This example cannot be faithfully described by a real-valued valuation
preserving the order of all of $\Gamma$. Sending $\Omega$ to a finite real
number would fail to preserve $\Omega>n$ for all $n$. The exact calculation
uses only the ordered-group comparisons and the support lemma.

The same threshold persists if $t^\Omega$ in the linear term is replaced by
\[
 c=t^\Omega\left(1+\sum_{j\ge2}t^{\,2-1/j}\right).
\]
This is a legitimate Hahn coefficient: its positive support is well ordered,
although infinitely many exponents lie below $2$ before the shift by
$\Omega$. The leading valuation remains $\Omega$, and the normalized shell
polynomial remains $1+X$. Thus the theorem is not restricted to coefficients
with discrete or finitely generated exponent support.

\section{Nonresonant Hahn linearization need not be coherent}\label{sec:noncoherent}

The extra conclusion in \cref{thm:coherent} cannot be imported into
\cref{lem:nonresonant}. The obstruction is ordinary small-divisor divergence
in the exponent-zero coefficient, not a failure of Hahn summability.

\begin{proposition}[A divergent ordinary linearizer survives on the monad]
\label{prop:noncoherent}
There exists an ordinary complex polynomial $f(X)=\alpha X+X^2$ with
$\alpha$ not a root of unity such that its normalized formal conjugacy is
strongly Hahn-evaluable and isometric on $\mm$ in every nontrivial workspace
$K$, but does not belong to $\C\{X\}((t^\Gamma))$. In particular, it has
no common-domain representative as a normalized formal conjugacy.
\end{proposition}

\begin{proof}
Choose $\alpha=e^{2\pi i\theta}$ with $\theta$ an irrational non-Brjuno
number. Yoccoz's quadratic theorem states that the germ
$\alpha X+X^2$ is analytically linearizable exactly when $\theta$ is Brjuno
\cite{yoccoz}. Thus its unique normalized formal linearizer
$H(X)\in\C[[X]]$ has zero ordinary radius of convergence.

All coefficients of this formal series lie in the constant field, so its
Hahn support is simply $\{0\}$. By \cref{lem:evaluation}, both $H$ and
its formal inverse are strongly evaluable and isometric on $\mm$.
The exact maximal valuation ball is $B_0$ by \cref{thm:exact}.
But the coefficient at Hahn exponent zero is the divergent ordinary series
$H$ itself. Therefore the series is not in the radius-free convergent-germ
ring, let alone in a common-domain ring.
\end{proof}

For an explicit choice, use continued-fraction convergents with denominators
$q_{-1}=0$, $q_0=1$, and choose partial quotients recursively by
\[
 a_{n+1}=2^{q_n^2},\qquad
 q_{n+1}=a_{n+1}q_n+q_{n-1},\qquad
 \theta=[0;a_1,a_2,\ldots].
\]
Then
\[
 \frac{\log q_{n+1}}{q_n}\ge q_n\log2,
\]
so the Brjuno sum $\sum_n(\log q_{n+1})/q_n$ diverges. This construction
specifies a multiplier; no numerical approximation to it is used.

\begin{lemma}[Faithfulness of infinitesimal evaluation]\label{lem:faithful}
If $0\ne A\in\C[[X]]((t^\Gamma))$ and $\Gamma\ne\{0\}$ is divisible,
then $A(u)\ne0$ for some $u\in\mm$. Evaluation on the whole monad is
therefore injective on this ring.
\end{lemma}

\begin{proof}
Let $\gamma_0$ be the least Hahn exponent of $A$, and let $n$ be the
least degree of the nonzero formal series $A_{\gamma_0}(X)$.
If another Hahn exponent occurs, let $\Delta>0$ be the gap from $\gamma_0$
to the next exponent. Choose $s>0$ with $ns<\Delta$; divisibility supplies
such an $s$ when $n>0$, and any positive $s$ works when $n=0$.
If no other exponent occurs, choose any $s>0$.

At $u=t^s$, the exponent $\gamma_0+ns$ appears with nonzero coefficient.
Higher degrees at $\gamma_0$ and all other Hahn exponents contribute only
larger exponents. Thus $A(u)\ne0$.
All evaluations are legitimate: after factoring out $t^{\gamma_0}$,
the remaining exponents lie in the monoid generated by a well-ordered
positive set, so \cref{lem:evaluation} applies.
\end{proof}

The lemma rules out an evasion by changing the representation: a different
common-domain family cannot represent the divergent conjugacy as the same
function on all infinitesimals. Likewise, a normalized common-domain
conjugacy satisfying the same functional identity would satisfy it as a
formal identity, and uniqueness of \eqref{eq:recurrence} would force the
same divergent exponent-zero coefficient.

There is no conflict with Yoccoz's theorem. An ordinary divergent series can
be Hahn-summable at an infinitesimal argument because each successive power
introduces a new positive exponent. This is a change of analytic category,
not a proof of classical analytic linearizability. Nor is every nonresonant
case noncoherent; the proposition establishes failure of a universal
common-domain assertion, not an if-and-only-if classification.

\section{Workspace invariance and the full surcomplex class}\label{sec:workspace}

\subsection{Enlarging a set-sized workspace}

\begin{theorem}[Workspace invariance]\label{thm:workspace}
Let $\Gamma\subseteq\Gamma'$ be divisible ordered abelian groups and view
$K=\C((t^\Gamma))$ as a subfield of
$K'=\C((t^{\Gamma'}))$. For a polynomial $f\in K[z]$ satisfying
\cref{thm:exact}, the threshold $r$, normalized formal conjugacy, and
resonant shell polynomial are unchanged by this extension.

The conjugacy and its inverse strongly evaluate at every point of
$B_r(K')$, including those absent from $K$. The finite shell root set,
algebraic multiplicities, exact periods, and multiplier formulas gain no new
points and do not change. The common-domain coefficient functions in the
resonant case are also unchanged.
\end{theorem}

\begin{proof}
The embeddings preserve valuation, residue, and all finite coefficient
arithmetic. Thus the return polynomial and the finite maximum defining $r$
are identical. The recurrence \eqref{eq:recurrence} gives the same formal
coefficients in both fields.

For a new point $z\in B_r(K')$, the normalized input $z/t^r$ has positive
well-ordered support in $\Gamma'$. Its union with the already constructed
positive support certificate is still well ordered. Hence the evaluation
and composition proof of \cref{lem:evaluation} applies without change.

Every relevant root is a root of a polynomial over the algebraically closed
field $K$. That polynomial already factors completely into linear factors
in $K[z]$, and the same finite factorization remains valid in $K'[z]$.
It therefore has no additional roots in $K'$, with no change of
multiplicity. All remaining assertions concern the same formal coefficients,
ordinary residue polynomial, and uniquely normalized differential recursion.
\end{proof}

\subsection{Applying the theorem to \texorpdfstring{$\No[i]$}{No[i]}}

Conway normal forms identify the surreal field with the appropriate
class-sized Hahn field of real coefficients; this viewpoint is recalled,
for example, in Bournez--Guilmant \cite{bournez-guilmant}.
Complexifying gives normal forms with complex coefficients. For a
set-sized divisible ordered subgroup $\Gamma\subseteq\No$, the substitution
\begin{equation}\label{eq:surreal-embedding}
 t^\gamma\longmapsto\omega^{-\gamma}
\end{equation}
embeds $\C((t^\Gamma))$ in $\SC=\No[i]$. The minus sign converts
increasing Hahn supports to decreasing surreal monomial supports.

Every set of surcomplex numbers lies in one such workspace: take the union
of the exponent supports of their normal forms, and then its divisible
additive span. This is a set. Adjoining one positive exponent ensures that
the workspace is nontrivial.

\begin{corollary}[Surcomplex polynomial dynamics]\label{cor:surcomplex}
Let $f\in\SC[z]$ be a nonlinear polynomial fixing zero, with unit
multiplier for the natural Hahn valuation, and with that multiplier not a
root of unity. Then the formulas and conclusions of
\cref{thm:exact,thm:shell,thm:coherent,thm:stability} hold on the corresponding
class-sized ball in $\SC$, interpreted pointwise by set-sized Hahn sums.
The threshold and finite shell cycles are independent of the workspace used
to calculate them.
\end{corollary}

\begin{proof}
Place the finitely many coefficients of $f$ in one divisible workspace.
Construct its countable sequence of formal conjugacy coefficients there.
For any particular argument, enlarge the workspace to include that argument
and apply \cref{thm:workspace}. Two choices agree in their common larger
workspace, so they define one class function pointwise. The finite
obstructing root set already lies in the initial algebraically closed
workspace. No sum in this construction is indexed by a proper class.
\end{proof}

The corollary may be interpreted in a class theory such as NBG, or with
definable classes in ZFC. It does not require a proper-class polynomial ring
as a set, and it does not replace Hahn summation by fine-topological
convergence. In particular, no Banach contraction principle, Cauchy-sequence
limit, or ordinary surcomplex path is used in any proof.

\section{Computation, proof dependencies, and limitations}\label{sec:verification}

\subsection{An exact finite procedure}

For polynomial inputs with effective coefficient arithmetic, the geometric
certificate can be calculated without constructing $H$.
\begin{enumerate}[label=\arabic*.]
\item Determine $\alpha=\bar\lambda$. If it is not a root of unity,
compute the finite maximum $\eta$ in \eqref{eq:eta} and return $r=\eta$.
\item If its order is $q$, compute the finite composition $g=f^{\circ q}$,
$c=\lambda^q-1$, and $\delta=v(c)$.
\item Compute every nonzero nonlinear return coefficient $b_n$, including
actual cancellations, and take the maximum in \eqref{eq:exact-r}.
\item Scale by $\rho=t^r$ and reduce \eqref{eq:Q-shell} to obtain $P$.
Its degree gives the shell multiplicity. Factoring $P$ determines residue
directions, and $XP'(X)$ gives the leading multiplier ratios at simple roots.
\end{enumerate}
This is a finite algebraic procedure relative to the required oracles, not a
claim of universal computability for arbitrary surcomplex inputs. In
particular, equality and valuation of arbitrary Hahn expressions need not be
effective; recognizing root-of-unity status for an arbitrary complex
coefficient requires a presentation allowing that decision. Moreover,
$d^q$ may be very large. The formula is an exact characterization and a
certificate, not a polynomial-time complexity theorem.

When coefficients are available only to a specified precision, the stability
theorem tells which strict inequalities certify that the computed first face
is already correct. At a tie or a possible cancellation, more precision may
be necessary. Treating a truncated zero as an exact zero would invalidate the
maximum.

\subsection{The accompanying finite checks}

The archive contains \texttt{verify.py} and the recorded output
\texttt{verification.txt}. The script uses exact symbolic and rational
arithmetic, not floating-point approximations. Its 54 checks cover:
\begin{itemize}
\item the full second iterate of the general cubic, the tuned factorization,
and the three valuation regimes with explicit integer-exponent
representatives;
\item the shell multipliers, the transition discriminant, and the
logarithmic derivatives of the three leading coherent conjugacies;
\item Schr\"oder identities through degree nine for several exact rational
and complex multipliers, and the exact binomial multiplier identity;
\item integrality and residue of the normalized tuned conjugacy through
degree nine, including the coefficients $[X^5]H_0=1/2$ and
$[X^9]H_0=5/8$.
\end{itemize}
The recorded run reports 54 of 54 checks passing. These tests catch arithmetic
and expansion errors in displayed examples. They do not establish the
arbitrary-rank support lemma, the transfinite lifting argument, or maximality
for every polynomial, and they are not an implementation of general Hahn
fields. The proofs, not the tests, carry those claims.

\subsection{The proof dependencies}

The argument uses two external structural facts: the positive-support lemma
and algebraic closure of $\C((t^\Gamma))$. Starting there, the dependence is
\[
\begin{gathered}
 \text{uniform support + triangular recurrence}\\
 \Downarrow\\
 \text{nonresonant and divisible-displacement lemmas}\\
 \Downarrow\\
 \text{exact ball + polynomial root obstruction}\\
 \Downarrow\\
 \text{shell symmetry, cycle count, multipliers, stability}.
\end{gathered}
\]
The common-domain theorem branches from the divisible-displacement lemma
and uses ordinary holomorphic integration to solve the transfinite family
\eqref{eq:ode-gamma}. The noncoherent example additionally uses Yoccoz's
classical quadratic theorem. Workspace invariance uses finite factorization
and the same support certificates, not a new analytic-completion theorem.

\subsection{Why equal characteristic zero matters}

The cancellation in \eqref{eq:denomcancel} leaves a factor $n-1$ in the
denominator. In this article every such nonzero integer is a valuation unit.
That property fails in mixed characteristic. One should therefore not
transport the exact finite-return formula or the shell-periodicity conclusion
to $\C_p$. Indeed, Lindahl gives $p$-adic quadratic-like examples whose
linearization bounding sphere contains no periodic point
\cite[Corollary~C]{lindahl-padic}. The contrast identifies a hypothesis in
our proof; it is not a contradiction with those examples.

\subsection{The claims that are deliberately not made}

The main theorem covers nonlinear polynomials in one variable, with unit
multiplier not a root of unity, over complex Hahn fields of arbitrary
ordered-group rank. It does not cover all infinite Hahn-analytic power
series: there a common support certificate, a finite or attained dominant
face, and analytic composition must be established separately. It does not
address actual root-of-unity multipliers, attracting or repelling multipliers,
or higher-dimensional simultaneous resonances.

Nor is $r$ asserted to describe every possible analytic continuation of the
conjugacy. It describes the maximal centered valuation ball. The common-domain
coefficient functions may supply additional local information in selected
residue directions, without yielding a larger centered valuation ball.
The shell is not an ordered-modulus circle, and is not a topological boundary.

The literature review was targeted. The strongest proposed contributions
are the exact resonant finite-return formula with its arbitrary-rank support
certificate, the shell-cycle and multiplier theorem, and the common-domain
lifting construction. The nonresonant rank-one radius, the existence of a
formal linearizer, and the underlying support and algebraic-closure machinery
are not claimed as new. The explicit cancellation family is a worked theorem
showing a strict gain over a coefficient-size bound, not evidence that every
near-resonant family has the same phase diagram.

Finally, this manuscript is an AI-assisted research draft. It has not been
refereed or formally verified. The mathematical statements are supplied with
proofs for expert assessment; the finite checks and targeted search should
not be mistaken for certification of either correctness or priority.

\section{Conclusion}

For indifferent surcomplex polynomial dynamics, the first nonlinear
obstruction to linearization can be detected by a finite return map. When
the residue multiplier has order $q$, the largest valuation among the nonzero
fixed points of $f^{\circ q}$ is exactly the threshold of the normalized
linearization ball. Those closest fixed points have exact period $q$, and a
finite complex initial polynomial records their directions, multiplicity,
and leading return multipliers.

The analytical content is not exhausted by that finite algebra. A uniform
positive-support certificate turns the formal conjugacy into a well-defined
function at every infinitesimal scale in an arbitrary-rank workspace.
Resonance then allows a further upgrade: the coefficient functions all solve
ordinary differential equations on one common disk. The non-Brjuno example
shows why this last step cannot be assumed merely from infinitesimal Hahn
evaluability.

The cubic transition makes the difference concrete. A cancellation invisible
to the basic coefficient scale changes the exact threshold from $\delta/2$
to $\delta/4$, creates a second closest two-cycle, and changes the leading
coherent coordinate from a square-root expression to a fourth-root
expression. The same proof works at genuinely non-Archimedean valuation
rank and survives passage to the full surcomplex class, without invoking a
proper-class sum or a nonexistent sequence limit.

\clearpage
\phantomsection
\addcontentsline{toc}{section}{References}
\begin{thebibliography}{9}

\bibitem{repo}
Vladimir Reshetnikov (repository maintainer),
\emph{Surreal}, research documentation.
Snapshot inspected: commit
\texttt{39f2be6667ade51bca2b45daa47e289d69c09764},
21 September 2026.
Catalogues inspected: \texttt{docs/README.md},
\texttt{docs/surcomplex/analysis/README.md}, and
\texttt{docs/surcomplex/analytic-geometry/README.md}.
The catalogues identify the reports as AI-assisted, unrefereed drafts.
\url{https://github.com/VladimirReshetnikov/Surreal/tree/39f2be6667ade51bca2b45daa47e289d69c09764/docs}.

\bibitem{neumann}
B.~H. Neumann,
On ordered division rings,
\emph{Transactions of the American Mathematical Society}
\textbf{66} (1949), 202--252.
\url{https://www.jstor.org/stable/1990552}.

\bibitem{poonen}
Bjorn Poonen,
Maximally complete fields,
\emph{L'Enseignement Math\'ematique} (2)
\textbf{39} (1993), 87--106.
In particular, Corollary~4.
\url{https://math.stanford.edu/~conrad/Perfseminar/refs/poonencomplete.pdf}.

\bibitem{lindahl-equal}
Karl-Olof Lindahl,
Linearization in ultrametric dynamics in fields of characteristic zero---equal
characteristic case,
\emph{p-Adic Numbers, Ultrametric Analysis and Applications}
\textbf{1} (2009), 307--316.
\href{https://doi.org/10.1134/S2070046609040049}{doi:10.1134/S2070046609040049}.
Author version: arXiv:1111.1993 (2011),
\url{https://arxiv.org/html/1111.1993v1}.

\bibitem{lindahl-padic}
Karl-Olof Lindahl,
The size of quadratic $p$-adic linearization disks,
arXiv:1311.4477 (2013), version~1.
In particular, Corollary~C.
\url{https://arxiv.org/html/1311.4477v1}.

\bibitem{yoccoz}
Jean-Christophe Yoccoz,
Th\'eor\`eme de Siegel, nombres de Bruno et polyn\^omes quadratiques,
in \emph{Petits diviseurs en dimension 1},
\emph{Ast\'erisque} \textbf{231} (1995), 1--88.
\url{https://www.numdam.org/book-part/AST_1995__231__1_0/}.

\bibitem{bournez-guilmant}
Olivier Bournez and Quentin Guilmant,
Surreal fields stable under exponential and logarithmic functions,
arXiv:2201.08199 (2022).
Cited for the recalled Hahn-series viewpoint on surreal normal forms, not as
an origin of Conway's normal-form theorem.
\url{https://arxiv.org/abs/2201.08199}.

\end{thebibliography}
\end{document}
